\chapter{Python code used for  Estimation of Dimensions}
\label{app:ship_dimension_generator}

The following Python program was developed to generate feasible
combinations of the principal dimensions of the proposed vessel.
The program evaluates combinations of length overall (LOA), breadth
(B), depth (D), block coefficient ($C_B$), and displacement, and
calculates the corresponding waterline length, breadth at the design
waterline, draft, and principal dimensional ratios.

\begin{lstlisting}[language=Python,
    caption={Python program for generating feasible ship main dimensions},
    label={lst:ship_dimension_generator}]
import csv

# ============================================================
# SHIP MAIN DIMENSION GENERATOR
# ============================================================

# MAIN DIMENSIONS
LOA_MIN = 114
LOA_MAX = 118
LOA_STEP = 1

B_MIN = 14
B_MAX = 16
B_STEP = 1

D_MIN = 9
D_MAX = 10
D_STEP = 1

# BLOCK COEFFICIENT
CB_VALUES = [0.380, 0.39, 0.395, 0.400]

# DISPLACEMENT
DISP_MIN = 3600
DISP_MAX = 3850
DISP_STEP = 50

# CONSTANTS
LWL_RATIO = 0.975
B_DWL_RATIO = 0.94
RHO = 1.025


# ============================================================
# GENERATE INTEGER VALUES
# ============================================================

LOA_VALUES = range(LOA_MIN, LOA_MAX + 1, LOA_STEP)
B_VALUES = range(B_MIN, B_MAX + 1, B_STEP)
D_VALUES = range(D_MIN, D_MAX + 1, D_STEP)
DISP_VALUES = range(DISP_MIN, DISP_MAX + 1, DISP_STEP)


# ============================================================
# CALCULATE POSSIBLE DIMENSIONS
# ============================================================

results = []

for LOA in LOA_VALUES:

    # LWL is calculated from LOA
    LWL = 0.975 * LOA

    for B in B_VALUES:

        # Breadth at DWL
        B_DWL = 0.94 * B

        for D in D_VALUES:

            for CB in CB_VALUES:

                for displacement in DISP_VALUES:

                    # ----------------------------------------
                    # DISPLACEMENT EQUATION
                    #
                    # Delta = rho x LWL x B_DWL x T x CB
                    #
                    # T = Delta / (rho x LWL x B_DWL x CB)
                    # ----------------------------------------

                    T = displacement / (
                        RHO *
                        LWL *
                        B_DWL *
                        CB
                    )

                    # Ratios
                    L_B = LOA / B
                    B_T = B / T
                    D_T = D / T

                    # Draft must be less than depth
                    if T < D:

                        results.append([
                            LOA,
                            LWL,
                            B,
                            B_DWL,
                            D,
                            CB,
                            displacement,
                            T,
                            L_B,
                            B_T,
                            D_T
                        ])


# ============================================================
# DISPLAY RESULTS
# ============================================================

print()
print("=" * 120)
print("POSSIBLE SHIP MAIN DIMENSIONS")
print("=" * 120)

print()

print(
    f"{'No.':<5}"
    f"{'LOA':<7}"
    f"{'LWL':<9}"
    f"{'B':<7}"
    f"{'B_DWL':<9}"
    f"{'D':<7}"
    f"{'CB':<8}"
    f"{'Disp':<9}"
    f"{'Draft':<9}"
    f"{'L/B':<8}"
    f"{'B/T':<8}"
    f"{'D/T':<8}"
)

print("-" * 120)


for i, row in enumerate(results, 1):

    (
        LOA,
        LWL,
        B,
        B_DWL,
        D,
        CB,
        displacement,
        T,
        L_B,
        B_T,
        D_T
    ) = row

    print(
        f"{i:<5}"
        f"{LOA:<7}"
        f"{LWL:<9.3f}"
        f"{B:<7}"
        f"{B_DWL:<9.2f}"
        f"{D:<7}"
        f"{CB:<8.3f}"
        f"{displacement:<9}"
        f"{T:<9.3f}"
        f"{L_B:<8.2f}"
        f"{B_T:<8.2f}"
        f"{D_T:<8.2f}"
    )


# ============================================================
# SAVE RESULTS TO CSV
# ============================================================

filename = "ship_possible_dimensions.csv"

with open(
    filename,
    "w",
    newline="",
    encoding="utf-8"
) as file:

    writer = csv.writer(file)

    writer.writerow([
        "LOA (m)",
        "LWL (m)",
        "Breadth B (m)",
        "Breadth at DWL (m)",
        "Depth D (m)",
        "CB",
        "Displacement (t)",
        "Draft T (m)",
        "L/B",
        "B/T",
        "D/T"
    ])

    for row in results:
        writer.writerow(row)


# ============================================================
# FINAL MESSAGE
# ============================================================

print()
print("=" * 120)
print("TOTAL FEASIBLE COMBINATIONS =", len(results))
print("=" * 120)
print()
print("Results saved as:", filename)
\end{lstlisting}
